\subsection{How much of this is the uniform prior?}
\label{sec:priors}

Condition~\eqref{eq:margin} is an all-teammate-profile condition: it never
mentions the prior. Only $D$ and $\Dnull$ do. So the threshold question is well
posed for any full-support prior $\pi$, and Lemma~\ref{lem:necessity} survives
verbatim, because its argument is about an edge derivative vanishing and uses
only that attaining $\Dnull(K)$ forces each profile's mass onto its nearest
centres.

What does \emph{not} survive is the step from that lemma to 2-domination.

\begin{proposition}[Necessity, in the form that holds for any prior]
\label{prop:necessity-general}
Let $\pi$ have full support and let the record be flat at $(n,K)$ with centre set
$C$. Then every $s \notin C$ has at least two centres at distance $d(s,C)$.
\end{proposition}

This is 2-domination only when $d(s,C) = 1$ for every $s$, that is, when the
$\pi$-optimal code is \emph{dominating}. Under the uniform prior with
$K \ge \gamma(\Qn)$ it is, which is what Theorem~\ref{thm:kflat} relies on. Under
a tilted prior it need not be: the optimal code concentrates on heavy profiles
and can leave a light profile at distance $2$.

\paragraph{A counterexample, and what it retracts.}
Take $n = 3$, a product Bernoulli$(0.7)$ prior, $K = 4$, and
$C = \{011, 101, 110, 111\}$. This $C$ is the unique $\pi$-optimal $4$-median set,
with $\Dnull(4) = 0.2430$. It is not dominating: $000$ lies at distance $2$ from
it. Spreading each non-centre uniformly over its nearest centres gives a record
with $D = 0.2430 = \Dnull(4)$ exactly and margin $\gamma = 1/3$. It is flat.

An earlier version of this paper claimed that for every tilted product prior the
threshold is $2^n - (n-1)$, and that the free compression ratio therefore
collapses to $1 + O(n2^{-n})$. \textbf{That claim was wrong and we withdraw it.}
The error was exactly the step above---reading
Proposition~\ref{prop:necessity-general} as 2-domination when the optimal code is
not dominating. Searching all optimal codes under the corrected condition gives

\begin{center}\small
\begin{tabular}{@{}lrrr@{}}
\toprule
 & uniform & tilted ($q = 0.7$) & withdrawn claim \\
\midrule
$n = 3$ & $4$ & $4$ & $6$ \\
$n = 4$ & $7$ & $5$ & $13$ \\
\bottomrule
\end{tabular}
\end{center}

so a tilted prior can require \emph{fewer} labels than the uniform one, not
vastly more. The direction of the withdrawn claim was backwards.

\paragraph{What we can say, and what is open.}
Theorem~\ref{thm:kflat} is unaffected: it is stated for the uniform prior, where
optimal codes are dominating, and both its halves were verified independently.
For general priors we have Proposition~\ref{prop:necessity-general} and the
matching construction in the two cases above, but not a proof that the
construction attains a positive margin for every optimal code satisfying the
condition. We therefore state the general-prior threshold as \textbf{open}, with
the corrected necessary condition recorded and the 2-domination characterisation
established only for the uniform case.

The honest reading of the uniform prior is correspondingly weaker than we first
claimed. It is not that symmetry is uniquely favourable; it is that symmetry
creates \emph{ties}---at $n=3$, $K=4$ the uniform prior has $62$ optimal codes
and the tilted prior has one---and ties give the designer freedom to pick an
optimal code that is also well-connected. Whether that freedom helps or hurts
depends on the prior, and we no longer claim a direction.
